\documentclass[10pt,a4paper]{amsart}
\usepackage[margin=19mm]{geometry}
\usepackage{amsmath,amssymb,mathtools}
\usepackage[scr=rsfs,cal=euler]{mathalfa}
\usepackage{xfrac}
\usepackage[unicode,hidelinks,bookmarks=false]{hyperref}
\newcommand{\ie}{i.e.\ }
\newcommand{\cf}{cf.\ }
\newcommand{\resp}{respectively\ }
\newcommand{\Subsection}{\subsection}
\providecommand{\nobreakdash}{\nobreak\hbox{-}}
\setlength{\emergencystretch}{2em}
\multlinegap0pt
\usepackage{amssymb}
\usepackage[shortlabels]{enumitem}
\usepackage{tikz}
\usepackage{xcolor}
\usepackage{float}
\usepackage[normalem]{ulem}
\newtheorem{lemma}{Lemma}
\def\Z{\mathbb{Z}}
\title{Unimodality of $q$-Fibonomial coefficients for small cases}
\author{Brendan B. Connelly, Ezekiel Ito, Thomas C. Martinez, Olha Shevchenko, Kacey Yang}
\date{Source: arXiv:2605.12822v1 (CC BY 4.0)}
\begin{document}
\maketitle

\noindent\emph{Source and licence.} Unimodality of $q$-Fibonomial coefficients for small cases. Version arXiv:2605.12822v1 (CC BY 4.0), distributed by arXiv under the Creative Commons Attribution 4.0 International licence, http://creativecommons.org/licenses/by/4.0/, as recorded in the arXiv licence metadata for this exact version and shown on the abstract page https://arxiv.org/abs/2605.12822v1. Full licence notice: \texttt{licenses/arxiv-2605.12822-CC-BY-4.0.txt}.\par\smallskip

\noindent\textit{Editorial scope.} Counted unit: the article's lemma lem:Ak-greater-than-Bk (source0.tex lines 1109--1120). The article's complete original proof of it is delivered with it (source0.tex lines 1122--1156, one \texttt{proof} environment of 35 lines). No mathematics is written, repaired or re-derived: the statement and the proof are the article's own lines, in the article's own notation, and no retained line is substituted or re-flowed.  No further source-local context is needed: every cross-reference of every retained line resolves inside the two delivered spans.  Nothing was omitted from any retained span, so the package carries no omission record.  No retained line contains a citation, so the wrapper carries no bibliography and no bibliography entry.  No retained line contains a figure environment, an image-inclusion command or drawing code, so no image is delivered and the package declares no asset dependency.  Editorial formatting only, in addition: an AMS \texttt{amsart} preamble that restates the article's own declarations and macros used by the retained lines, namely the environments \texttt{lemma} on separate counters, together with the standard packages those lines need.  The retained environments print in this document as Lemma 1 in order of appearance, whereas the article prints the same items with the numbering of its own full document.  The complete native source of the article as distributed in the arXiv e-print for this exact version is bundled unchanged as \texttt{sources/200456/source0.tex}.
\begin{lemma}
    Suppose $a,b,c\in\mathbb{Z}$ with $a \leq b$ and $a,b$ odd. Define
    $$A(k) = \# \left \{ \ell \in [0, c - 1] \, \Bigm| \, \frac{k - a + 2}{2} \le \ell  \le \frac{k + 1}{2}, \; \ell \in \Z_{\geq 0} \right \},$$
    and 
    $$B(k) =  \# \left \{ \ell \in [0, c - 1] \, \Bigm|\, \frac{k - a - b + 2}{2} \le \ell \le \frac{k - b + 1}{2}, \; \ell \in \Z_{\geq0} \right \}.$$
    Then:
    \begin{enumerate}
        \item If $2c \leq a+b$, then $A(k) \geq B(k)$ for every $k \leq c + \frac{a+b}{2} - 3$.
        \item If $2c > a+b$, then $A(a+b-2) < B(a+b-2)$.
    \end{enumerate}
    \label{lem:Ak-greater-than-Bk}
\end{lemma}

\begin{proof}
Assume $2c \leq a+b$ and fix $k \leq c + \tfrac{a+b}{2} - 3$.
 
Since $2c \leq a+b$, we have $k \leq a+b-3$, so $\tfrac{k-a-b+2}{2} \leq -\tfrac{1}{2}$ and the lower bound on $\ell$ in $B(k)$ reduces to $0 \leq \ell$. Since $a \leq b$, we have $k \leq c + b - 3$, so $\tfrac{k-b+1}{2} \leq c - 1$ and the upper bound reduces to $\ell \leq \tfrac{k-b+1}{2}$. Hence
\begin{align}
    B(k)&=\# \left \{ \ell \in \Z \Bigm| 0 \le \ell \le \frac{k - b + 1}{2}\right \} \nonumber\\
    &=\begin{cases} 0, & k < b-1, \\ \left\lfloor \dfrac{k-b+1}{2} \right\rfloor + 1, & b-1 \leq k \leq c + \dfrac{a+b}{2} - 3. \label{eq:Bk-formula} \end{cases}
\end{align}
For $k < b-1$, then, $B(k) = 0 \leq A(k)$.
 
Now assume $k \geq b-1$. Then $\tfrac{k-a+2}{2} \geq \tfrac{1}{2}$, so the lower bound on $\ell$ in $A(k)$ reduces to $\tfrac{k-a+2}{2} \leq \ell$, while the upper bound is $\min\bigl(\tfrac{k+1}{2}, c-1\bigr)$, with $\tfrac{k+1}{2} \leq c-1$ equivalent to $k \leq 2c-3$. Since $a$ is odd, both endpoints $\tfrac{k-a+2}{2}$ and $\tfrac{k+1}{2}$ are integers when $k$ is odd and half-integers when $k$ is even. Thus
\begin{align}
A(k) &= \# \left \{ \ell \in \Z \, \Bigm|  \frac{k - a + 2}{2} \le \ell  \le \min \left(\frac{k + 1}{2},c-1\right) \right \} \nonumber \\
&= \begin{cases} 
    \dfrac{a-1}{2}, & b-1 \leq k \leq 2c-3,\ k \text{ even}, \\[5pt] 
    \dfrac{a+1}{2}, & b-1 \leq k \leq 2c-3,\ k \text{ odd}, \\[5pt]
    \left\lfloor \dfrac{a+2c-k-2}{2} \right\rfloor, & 2c-2 \leq k \leq c + \dfrac{a+b}{2} - 3. \label{eq:Ak-formula} 
\end{cases}
\end{align}
 
For $b-1 \leq k \leq 2c-3$, we have $k - b + 1 \leq a - 2$, so $B(k) \leq \lfloor \tfrac{a-2}{2} \rfloor + 1 = \tfrac{a-1}{2} \leq A(k)$, where the final equality uses $a$ odd.
 
For $2c-2 \leq k \leq c + \tfrac{a+b}{2} - 3$, $A$ is non-increasing and $B$ is non-decreasing, so it suffices to verify $A(c + \tfrac{a+b}{2} - 3) \geq B(c + \tfrac{a+b}{2} - 3)$. Then \eqref{eq:Bk-formula} and \eqref{eq:Ak-formula} give
\[
    A(c + \tfrac{a+b}{2} - 3) = \left\lfloor \tfrac{c + \tfrac{a-b}{2}+1}{2} \right\rfloor \geq \left\lfloor \tfrac{c + \tfrac{a-b}{2}}{2} \right\rfloor = B(c + \tfrac{a+b}{2} - 3).
\]
This proves~(1).

For~(2), suppose $2c > a+b$. We have 
$$B(a+b-2)= \# \left \{ \ell \in [0, c - 1] \, \Bigm|\, 0 \le \ell \le \frac{a -1}{2}, \; \ell \in \Z_{\geq 0} \right \}=\frac{a+1}{2}$$
since $a$ is odd. 
$$A(a+b-2)=\# \left \{ \ell \in [0, c - 1] \, \Bigm| \, \frac{b}{2} \le \ell  \le \frac{a+b-1}{2}, \; \ell \in \Z_{ \geq 0} \right \}=\frac{a-1}{2}$$
since $b, a+b-1$ are odd and $\frac{a+b-1}{2} < c-1$. We have $A(a+b-2)<B(a+b-2)$.

\end{proof}

\end{document}
